(1) 理想轴角 $\alpha$ 的Jones透光投影为 $P_\alpha=\mathbf e_\alpha\mathbf e_\alpha^\mathsf T$，$\mathbf e_\alpha=(\cos\alpha,\sin\alpha)^\mathsf T$。相邻轴角差 $\delta=\Theta/N$，每一次投影保留的场振幅因子是 $\cos\delta$，强度因子为 $\cos^2\delta$。

(2) 在入射初偏振已经归一为轴角0的条件下，最终出射Jones场为 $E_0\cos^N(\Theta/N)\mathbf e_\Theta$，所以 $\boxed{I_N=I_0\cos^{2N}(\Theta/N)}$。最终偏振方向确为 $\Theta$；若直接使用末端角度的一片理想片，强度为 $I_{\rm single}=I_0\cos^2\Theta$。

(3) 对 $N\gg1$，$\ln(I_N/I_0)=2N\ln\cos(\Theta/N)=-\Theta^2/N-\Theta^4/(6N^3)+\cdots$。因此 $\boxed{I_N/I_0=1-\Theta^2/N+O(N^{-2})\to1}$。更密集的缓慢投影序列能在偏振方向显著转动时仍接近无损，这是一种空间版Zeno极限。

(4) 经典说法是每片只滤除与局部轴正交的极小分量，损耗每步为 $O(\delta^2)$，总和 $N\delta^2\to0$；单光子说法是多次极近轴的偏振投影使每次生存概率趋1。输出强度始终不超过 $I_0$，不存在从偏振片取到额外光能。
